\documentclass[a4paper,addpoints]{article}
\usepackage[utf8]{inputenc}
\usepackage[italian]{babel}
\usepackage{amsfonts,amssymb,amsmath, amsthm}
\usepackage{xcolor}
\usepackage{pgf,tikz,pgfplots}
\pgfplotsset{compat=1.15}
\usepgfplotslibrary{fillbetween}
\usepackage{mathrsfs}
\usetikzlibrary{arrows}
\usetikzlibrary{calc}
\usepackage{tabu}

\usepackage{tikz}
\usetikzlibrary {shapes , arrows , fit , calc, positioning}
\tikzset{box/.style={draw, rectangle, rounded corners, thick, node distance =7 em, text width =6 em, text centered, minimum height =2.5 em}}
\tikzset{line/.style={draw, thick, -latex'}}

\usepackage{tikz-uml}
\tikzumlset{actor below=0.4cm, fill usecase=white, fill class=white, text=black, draw=black}

\begin{document}

\section*{Lostly --- Diagramma dei casi d'uso}

Il diagramma descrive le funzionalit\`a di Lostly dal punto di vista degli attori.
L'\emph{Utente} (registrato) pu\`o registrarsi e accedere, inserire un oggetto smarrito o ritrovato,
cercare gli oggetti e consultarne i dettagli, segnalare una possibile corrispondenza
(che genera la notifica all'altro utente), contattare l'altro utente per la restituzione
e gestire le proprie segnalazioni, aggiornandone lo stato.
L'\emph{Amministratore} eredita tutte le funzionalit\`a dell'Utente e in pi\`u gestisce
utenti e segnalazioni.

\begin{center}
\resizebox{\textwidth}{!}{%
\begin{tikzpicture}
    \begin{umlsystem}[x=4]{Lostly}
        % colonna sinistra
        \umlusecase[x=0, y=0,    width=3.6cm]{Registrazione e accesso}
        \umlusecase[x=0, y=-2.5, width=3.6cm]{Inserimento oggetto smarrito}
        \umlusecase[x=0, y=-5,   width=3.6cm]{Inserimento oggetto ritrovato}
        \umlusecase[x=0, y=-7.5, width=3.6cm]{Modifica e gestione delle proprie segnalazioni}
        \umlusecase[x=0, y=-10,  width=3.6cm]{Aggiornamento stato dell'oggetto}
        % colonna destra
        \umlusecase[x=6, y=0,    width=3.6cm]{Ricerca oggetti smarriti e ritrovati}
        \umlusecase[x=6, y=-2.5, width=3.6cm]{Visualizzazione dettagli oggetto}
        \umlusecase[x=6, y=-5,   width=3.6cm]{Segnalazione possibile corrispondenza}
        \umlusecase[x=6, y=-7.5, width=3.6cm]{Notifica possibili corrispondenze}
        \umlusecase[x=6, y=-10,  width=3.6cm]{Contatto tra utenti per la restituzione}
        % amministratore
        \umlusecase[x=3, y=-13, width=4.2cm]{Gestione utenti e segnalazioni}
    \end{umlsystem}

    \umlactor[x=0, y=-5]{Utente}
    \umlactor[x=0, y=-13]{Amministratore}

    \umlinherit{Amministratore}{Utente}

    % associazioni Utente
    \umlassoc{Utente}{usecase-1}
    \umlassoc{Utente}{usecase-2}
    \umlassoc{Utente}{usecase-3}
    \umlassoc{Utente}{usecase-4}
    \umlassoc{Utente}{usecase-6}
    \umlassoc{Utente}{usecase-8}
    \umlassoc{Utente}{usecase-10}

    % associazione Amministratore
    \umlassoc{Amministratore}{usecase-11}

    % relazioni tra casi d'uso
    \umlextend{usecase-5}{usecase-4}
    \umlextend{usecase-7}{usecase-6}
    \umlextend{usecase-8}{usecase-7}
    \umlextend{usecase-10}{usecase-7}
    \umlinclude{usecase-8}{usecase-9}
\end{tikzpicture}}
\end{center}

\end{document}
